% Generated draft: requires mathematical and source review.
\documentclass[11pt]{article}
\usepackage[margin=25mm]{geometry}
\usepackage{fontspec}
\setmainfont{Latin Modern Roman}
\usepackage{amsmath,amssymb,mathtools,unicode-math}
\setmathfont{Latin Modern Math}
\usepackage{enumitem,needspace,xurl,hyperref}
\setlength{\parindent}{0pt}
\setlength{\parskip}{5pt}
\setlength{\emergencystretch}{3em}
\newcommand{\N}{\mathbb{N}}
\newcommand{\Z}{\mathbb{Z}}
\newcommand{\Q}{\mathbb{Q}}
\newcommand{\R}{\mathbb{R}}
\newcommand{\C}{\mathbb{C}}
\newcommand{\F}{\mathbb{F}}
\newcommand{\A}{\mathbb{A}}
\newcommand{\PP}{\mathbb{P}}
\newcommand{\E}{\mathbb{E}}
\newcommand{\eps}{\varepsilon}
\newcommand{\dd}{\,\mathrm{d}}
\begin{document}
\section*{UnsolvedMath ID: 30000380}
\textbf{Source problem code:} catalogue entry
\subsection{Definitions and mathematical statement}
Let $\omega_1$ be the first uncountable ordinal and $[\omega_1]^2$ its set of unordered two-element subsets. A coloring $c:[\omega_1]^2\to C$ is $2$-bounded if every color occurs on at most two pairs:
\[
 \forall\gamma\in C,\qquad|c^{-1}(\{\gamma\})|\le2.
\]
A subset $H\subset\omega_1$ is polychromatic, or rainbow, when $c|_{[H]^2}$ is injective. This is a global bound on color fibers, not a bound at each vertex.

The source asks whether
\[
 \omega_1\rightarrow^*(\omega_1)^2_{2\text{-bdd}},
\]
meaning
\[
 \forall c\text{ that is }2\text{-bounded}\ \exists H\subset\omega_1,
 \quad\operatorname{otp}(H)=\omega_1\text{ and }c|_{[H]^2}\text{ injective}.
\]
For subsets of $\omega_1$, the order-type requirement is equivalent to uncountability. The negation would be a $2$-bounded coloring such that every uncountable $H$ contains two distinct pairs receiving the same color.

The literature assessment checked on 21 August 2026 reports a positive consistency result under strong forcing axioms and a negative instance under the continuum hypothesis, $2^{\aleph_0}=\aleph_1$, while saying no ZFC resolution was found. The precise forcing axioms and consistency hypotheses are not supplied. If both reported results apply to exactly the relation above and the two respective extensions of ZFC are consistent, they would yield relative independence rather than a single unconditional positive or negative ZFC assertion. Thus the record's open label cannot be treated as an independently verified claim that no independence result is known. The partition relation itself is preserved, with those supplied logical-status qualifications left explicit.
\subsection{Short English statement}
Must every coloring of pairs of countable ordinals, with each color used at most twice, have an uncountable rainbow subset? The supplied assessment reports contrasting forcing-axiom and CH behavior, so the underlying axioms and relative-consistency scope must accompany any status claim.
\subsection{Sources}
\begin{itemize}
\item \emph{Set Theory} (2006), Oberwolfach Reports, DOI 10.4171/owr/2005/55, pp. 3121--3174, \url{https://doi.org/10.4171/owr/2005/55}; proposed year supplied as 2005.
\item Literature assessment checked 21 August 2026: original report, \url{https://doi.org/10.4171/owr/2005/55}; ``indestructible colourings and rainbow Ramsey theorems,'' \url{https://arxiv.org/abs/0804.4548}.
\item Local UnsolvedMath registry, permanent ID $30000380$, exact source code \texttt{catalogue entry}, created 25 August 2026 and updated 16 September 2026. Corrected-and-verified status preserves the partition relation and its stated coloring equivalence.
\end{itemize}
\end{document}
